\documentclass[UTF8, a4paper, 12pt]{article}

% 中文支持
\usepackage[UTF8]{ctex}
\usepackage{zhnumber}

% 数学和科学符号
\usepackage{amsmath,amssymb,amsfonts}
\usepackage{mathtools}
\usepackage{mathrsfs}

% 图形和排版
\usepackage{graphicx}
\usepackage{booktabs,tabularx,multirow,multicol}
\usepackage{subfig}
\usepackage{wrapfig}
\usepackage{rotating}

% 超链接和颜色
\usepackage{hyperref}
\usepackage{xcolor}
\hypersetup{
    colorlinks=true,
    linkcolor=blue,
    citecolor=blue,
    urlcolor=blue
}

% 页面布局
\usepackage{geometry}
\geometry{a4paper, margin=2cm}

% 代码环境
\usepackage{listings}
\lstset{
    basicstyle=\ttfamily,
    numbers=left,
    frame=single,
    breaklines=true
}

% 自定义命令 - 术语定义
\def\LLM{LLM}
\def\LLMs{LLM}
\def\RoCo{RoCo}
\def\RoCoBench{RoCoBench}
\def\IK{IK}
\def\RRT{RRT}
\def\RL{RL}
\def\Grasp{Grasp}
\def\DoF{DoF}
\def\TAMP{TAMP}
\def\VLM{VLM}

% 旋转命令定义
\def\rot#1{\rotatebox{90}{#1}}

% 图表路径
\graphicspath{{./images/}}

\begin{document}

% 标题页
\begin{titlepage}
    \centering
    \vspace*{2cm}
    
    {\LARGE\bfseries\sffamily
    RoCo: 基于大语言模型的多机器人\\
    辩证协作\\
    \vspace{1cm}
    Dialectic Multi-Robot Collaboration\\
    with Large Language Models\\
    }
    
    \vspace{2cm}
    
    \Large
    赵曼迪\textsuperscript{1} \quad 施里亚·贾恩\textsuperscript{1} \quad 宋书然\textsuperscript{1}\\
    \vspace{0.5cm}
    \textbf{哥伦比亚大学}
    
    \vfill
    
    \normalsize
    \today
    
\end{titlepage}

% 摘要
\begin{abstract}
\textbf{摘要：}本文提出了一种利用预训练大语言模型（Large Language Model, LLM）进行高层通信和底层路径规划的多机器人协作新方法。机器人配备大语言模型，用于讨论和集体推理任务策略。随后，它们生成子任务计划和任务空间路点路径，由多臂运动规划器加速轨迹规划。我们还提供来自环境的反馈（如碰撞检测），并提示大语言模型智能体在上下文中改进其计划和路点。在评估方面，我们引入了 RoCoBench，这是一个包含6个任务的多机器人协作基准测试，涵盖了广泛的多机器人协作场景，并附带了用于智能体表征和推理的纯文本数据集。实验证明了我们方法的有效性——它在 RoCoBench 的所有任务中都取得了较高的成功率，并能适应任务语义的变化。我们的对话设置具有高度可解释性和灵活性——在真实世界实验中，我们展示了 RoCo 易于集成人类在环，用户可以与机器人智能体交流协作，共同完成任务。详见项目网站 project-roco.github.io 获取视频和代码。

\textbf{关键词：}多机器人协作，大语言模型，运动规划，自然语言交互，智能体对话

\end{abstract}

\newpage

% 目录
\tableofcontents
\newpage

% 1. 引言
\section{引言}

多机器人系统因其提升任务效率的潜力而备受关注，但同时也面临着诸多挑战。要使机器人有效地分工协作，需要对任务有高层级的理解，并考虑每个机器人的能力（如可达范围或有效载荷）。另一个挑战在于底层运动规划：随着机器人数量增加，配置空间呈指数级增长，寻找无碰撞的运动规划变得极其困难。此外，传统的多机器人系统通常需要针对特定任务进行工程设计，因此牺牲了泛化能力——由于大部分任务结构是预定义的，这些系统无法适应新场景或任务变化。

我们提出 RoCo（Robot Collaboration 的缩写），一种零样本多机器人协作方法，用于解决上述挑战。我们的方法包含三个关键组成部分：

\begin{itemize}
    \item \textbf{对话式任务协调：}为促进信息交换和任务推理，我们让机器人通过将每个机器人委托给大语言模型智能体进行对话来"相互交流"。这使得机器人能够以自然语言讨论任务，并具有高度的可解释性，便于监督。
    
    \item \textbf{由大语言模型生成的反馈改进型子任务计划：}多智能体对话以每个智能体的子任务计划（如拾取物体）结束。我们提供一组环境验证和反馈（如逆运动学失败或碰撞），直到大语言模型智能体提出有效的计划。
    
    \item \textbf{大语言模型引导的联合空间运动规划：}从经过验证的子任务计划中，我们提取机器人关节空间中的目标构型，并使用集中式 Rapidly-exploring Random Tree (RRT) 采样器来规划运动轨迹。我们探索了大语言模型一个较少被研究的能力：三维空间推理。给定任务空间中的起点、终点和障碍物位置，我们展示了大语言模型能够生成结合高层任务语义和环境约束的路点路径，并显著降低运动规划器的采样复杂度。
\end{itemize}

接下来，我们介绍 RoCoBench，这是一个包含6个多机器人操作任务的基准测试。实验证明，RoCo 在基准任务上是有效的：通过利用大语言模型捕获的常识知识，RoCo 能够灵活处理各种协作场景，无需任何针对任务的训练。

总而言之，我们提出了一种多机器人协作的新方法，并得到两个技术贡献的支持：1）一个基于大语言模型的多机器人框架（RoCo），该框架在处理大量任务时具有灵活性，改进了任务级协调和动作级运动规划；2）一个新的多机器人操作基准测试（RoCoBench），用于系统评估这些能力。它包括一套任务，旨在检查算法在处理不同任务语义（如顺序或并发）、不同工作空间重叠程度以及不同智能体能力（如可达范围和末端执行器类型）和 embodiment（如6 Degrees of Freedom (DoF) UR5、7 DoF Franka、20 DoF 人形机器人）时的灵活性和泛化能力。

% 2. 预备知识
\section{预备知识}

\subsection{任务假设}

我们考虑一个具有 $N$ 个机器人、有限时间范围 $T$ 和完整观测空间 $O$ 的协作多智能体任务环境。每个智能体 $n$ 有观测空间 $\Omega_n \subset O$。智能体可能具有非对称的观测空间和能力，这强调了通信的必要性。我们手动定义描述函数 $f$，将任务语义和观测在时间步 $t$ 转换为自然语言提示：$l_n^t = f^n(o_t)$，其中 $o_t \in \Omega_n$。我们还定义解析函数，将大语言模型输出（如文本字符串"PICK object"）映射到相应的子任务，该子任务可以由一个或多个夹爪目标构型来描述。

\subsection{多臂路径规划}

给定一个子任务，我们使用集中式运动规划来为所有机器人找到到达其目标构型的轨迹。设 $X \in \mathbb{R}^d$ 表示所有 $N$ 个机械臂的联合配置空间，$X_{ob}$ 为配置空间中的障碍物区域，则无碰撞空间 $X_{free} = X \setminus X_{ob}$。给定初始条件 $x_{init} \in X_{free}$ 和目标区域 $X_{goal} \in X_{free}$，运动规划器找到满足 $\sigma^*(0) = x_{init}$，$\sigma^*(1) \in X_{goal}$ 的最优 $\sigma^* : [0, 1] \rightarrow X$。得到的 $\sigma^*$ 由机器人的机械臂关节控制器用于开环执行。

% 3. 基于大语言模型的多机器人协作
\section{基于大语言模型的多机器人协作}

我们提出 RoCo，这是一种利用大语言模型进行机器人通信和运动规划的多机器人协作新方法。我们方法中的三个关键组成部分如图2所示，具体描述如下。

\subsection{基于大语言模型的多智能体对话}

我们假设多智能体任务环境具有非对称观测空间和技能能力，这意味着智能体在没有相互通信的情况下无法进行有意义的协调。我们利用预训练的大语言模型来促进这种通信。具体来说，在每次环境交互之前，我们设置一轮对话，其中每个机器人被委托给一个大语言模型生成的智能体，该智能体接收该机器人独有的信息，并且必须严格根据其角色（如"我是 Alice，我能到达……"）进行回应。

对于每个智能体的大语言模型提示，我们使用共享的总体结构，但包含智能体特定的内容，根据每个机器人的个人状态进行调整。提示由以下关键组成部分：

\begin{enumerate}
    \item \textbf{任务上下文（Task Context）：}描述任务的总体目标。
    \item \textbf{轮次历史（Round History）：}来自先前轮次的对话和执行的动作。
    \item \textbf{智能体能力（Agent Capability）：}智能体可用的技能和约束。
    \item \textbf{通信指令（Communication Instruction）：}如何回应其他智能体并正确格式化输出。
    \item \textbf{当前观测（Current Observation）：}每个智能体状态的独特信息，加上当前对话中的先前回复。
    \item \textbf{计划反馈（Plan Feedback）：}（可选）先前子任务计划失败的原因。
\end{enumerate}

我们使用以下通信协议来推动对话前进：每个智能体被要求（在提示的通信指令部分给出指令）在回复结束时在两个选项之间做出决定：1）指示其他人继续讨论；2）总结每个人的动作并提出最终动作提案；只有当每个智能体在当前对话轮次中至少回复一次时，才允许第二个选项。此协议旨在允许智能体自由讨论，同时保证在有限的交换次数内提出一个子任务计划。

\subsection{大语言模型生成的子任务计划}

当一轮对话结束时，最后发言的智能体会总结一个"子任务计划"，其中每个智能体获得一个子任务（如拾取物体），并可选地获得任务空间中的一系列三维路点。这个子任务计划首先通过一组验证，然后才进入执行。如果任何检查失败，反馈会被附加到每个智能体的提示中，并开始另一轮对话。验证按以下顺序进行，每个验证假设前一个检查已通过（例如，计划必须先被解析才能检查任务和智能体约束）：

\begin{enumerate}
    \item \textbf{文本解析（Text Parsing）：}确保计划遵循所需格式并包含所有必需关键词
    \item \textbf{任务约束（Task Constraints）：}检查每个动作是否符合任务和智能体约束
    \item \textbf{逆运动学（IK）：}检查每个机械臂的目标姿态是否可通过逆运动学实现
    \item \textbf{碰撞检测（Collision Checking）：}检查逆运动学求解的机械臂关节构型是否会导致碰撞
    \item \textbf{有效路点（Valid Waypoints）：}可选地，如果任务需要路径规划，每个中间路点必须是可逆运动学求解且无碰撞的，且所有步骤应均匀分布
\end{enumerate}

允许智能体重新规划，直到达到最大尝试次数，之后当前轮次结束而不执行任何动作，下一轮开始。如果任务在有限的轮次数内未完成，则认为该回合失败。

\subsection{大语言模型引导的联合空间运动规划}

一旦子任务计划通过所有验证，我们将其与逆运动学结合，为所有机械臂产生一个联合目标构型，并可选地，任务空间路点路径的每一步产生一个中间目标构型。目标构型被传递给基于 RRT 的多臂运动规划器，该规划器跨所有机械臂联合规划，并为每个机器人输出运动轨迹以在环境中执行，然后任务进入下一轮。

\textbf{玩具示例：大语言模型的三维空间推理能力。}作为一个激励性的示例，我们让一个大语言模型（GPT-4）在三维网格中规划多智能体路径。我们随机采样三个智能体的（起点，终点）坐标和一组障碍物，并提示 GPT-4 规划无碰撞路径。在对失败计划提供反馈的情况下，允许最多5次尝试，GPT-4 在30次运行中获得86.7\%的成功率，平均使用2.73次尝试。如图3所示为示例输出。详见附录13了解更多细节和结果。虽然这种能力令人鼓舞，但我们发现随着网格大小和障碍物数量的增加，它也存在局限性。

% 4. 基准测试
\section{基准测试}

RoCoBench 是一套包含6个多机器人协作任务的基准测试，场景设置为桌面操作。这些任务涉及对大语言模型来说语义上易于理解的常识物体，并涵盖需要不同机器人通信和协调行为的协作场景。详见附录10了解基准测试的详细文档。我们强调三个定义每个任务的关键属性，总结于表1：

\begin{enumerate}
    \item \textbf{任务分解：}任务是否可以分解为可以并行完成或按一定顺序完成的子部分。RoCoBench 中有三个任务具有顺序性质（如"制作三明治"任务要求按正确顺序堆叠食物），而其他三个任务可以并行执行（如"包装杂货"任务中的物体可以按任意顺序放入箱子）。
    
    \item \textbf{观测空间：}每个机器人智能体接收多少任务和环境信息。有三个任务提供共享的任务工作空间观测，而其他三个任务设置更加非对称，机器人必须相互询问以交换知识。
    
    \item \textbf{工作空间重叠：}运行中的机器人之间的接近程度；我们将每个任务从低、中或高进行排序，其中更高的重叠需要更仔细的低层协调（如"移动绳子"任务需要共同操控同一个物体）。
\end{enumerate}

\begin{table}[htbp]
\centering
\caption{RoCoBench 任务关键属性概览}
\begin{tabular}{lcccccc}
\toprule
任务 & 分类方块 & 清扫地面 & 包装杂货 & 整理橱柜 & 移动绳子 & 制作三明治 \\
\midrule
任务分解 & 顺序 & 顺序 & 并行 & 并行 & 并行 & 顺序 \\
观测空间 & 共享 & 共享 & 非对称 & 非对称 & 非对称 & 共享 \\
工作空间重叠 & 低 & 低 & 中等 & 中等 & 高 & 高 \\
\bottomrule
\end{tabular}
\end{table}

% 5. 实验
\section{实验}

\textbf{概述。}我们使用 RoCoBench 设计了一系列实验来验证我们的方法。在第5.1节，我们评估 RoCo 相对于不使用对话的 oracle 大语言模型规划器的任务性能，并对 RoCo 中对话提示的不同组成部分进行消融实验。第5.2节通过实验证明了大语言模型提出的三维路点在多臂运动规划中的好处。第5.3节包含定性结果，展示了 RoCo 的灵活性和适应性。其他实验结果（如失败分析）见附录12。

\subsection{RoCoBench 主要结果}

\textbf{实验设置。}我们使用 GPT-4 进行所有主要结果实验。除了我们的方法"对话"（Dialog）外，我们还设置了 oracle 大语言模型规划器"集中计划"（Central Plan），它接收完整的环境观测、所有机器人能力的信息和相同的计划反馈，并提示大语言模型一次为所有机器人规划动作。我们还对 RoCo 的提示组成部分进行两种消融实验：第一种移除对话和过去轮次的动作历史，即"无历史对话"（Dialog w/o History）。第二种移除环境反馈，即"无反馈对话"（Dialog w/o Feedback），其中失败的动作计划被丢弃，智能体被提示继续讨论而不提供详细失败原因。为弥补缺乏重新规划轮次，每个回合给予两倍的事件长度预算。

\textbf{评估指标。}在每个事件有限的轮次和每轮最大重新规划尝试次数的条件下，我们从三个指标评估任务性能：1）有限轮次内的任务成功率；2）智能体成功完成一个事件所需的环境步骤数，这衡量了机器人策略的效率；3）每个回合执行环境动作前平均使用的重新规划尝试次数——这反映了智能体理解和利用环境反馈来改进计划的能力。总体而言，如果任务成功率更高、需要更少的环境步骤且需要更少的重新规划次数，则方法被认为更好。

\textbf{结果。}评估结果如表2所示。我们注意到，尽管接收到的信息较少，对话智能体有时能取得与 oracle 规划器相当的性能。特别是在"分类方块"任务中，智能体能够通过对话找到相互帮助的策略，但 oracle 在尝试同时满足所有智能体约束时犯了错误。虽然移除历史信息或计划反馈轮次不会对某些任务的性能产生负面影响，但包含两者的完整提示取得了最好的整体结果。最后，在"包装杂货"任务中，oracle 规划器在路点规划方面表现出更强的能力，展示了更好地整合反馈并在单个坐标步骤上改进的能力。

\begin{table}[htbp]
\centering
\caption{RoCoBench 评估结果。我们报告了每个任务20次运行的平均成功率（$\uparrow$）、成功运行中的平均步骤数（$\downarrow$）以及所有运行的平均重新规划尝试次数（$\downarrow$）。}
\begin{tabular}{lcccccc}
\toprule
方法 & 分类方块 & 清扫地面 & 包装杂货 & 整理橱柜 & 移动绳子 & 制作三明治 \\
\midrule
Oracle 规划器 &  &  &  &  &  &  \\
成功率 & $0.50 \pm 0.11$ & $0.70 \pm 0.10$ & $0.82 \pm 0.06$ & $0.90 \pm 0.07$ & $1.00 \pm 0.00$ & $0.96 \pm 0.04$ \\
步骤，重新规划 & $2.3, 3.9$ & $8.6, 2.6$ & $11.1, 3.9$ & $4.0, 2.7$ & $8.4, 2.0$ & $8.8, 1.2$ \\
\midrule
无历史对话 &  &  &  &  &  &  \\
成功率 & $0.65 \pm 0.11$ & $0.73 \pm 0.11$ & $0.48 \pm 0.11$ & $1.00 \pm 0.00$ & $0.00 \pm 0.00$ & $0.33 \pm 0.12$ \\
步骤，重新规划 & $3.7, 3.1$ & $5.8, 1.4$ & $9.2, 3.1$ & $4.2, 1.4$ & N/A, 1.0 & $9.6, 1.8$ \\
\midrule
无反馈对话 &  &  &  &  &  &  \\
成功率 & $0.45 \pm 0.11$ & $0.53 \pm 0.13$ & $0.35 \pm 0.10$ & $0.70 \pm 0.10$ & $0.95 \pm 0.05$ & $0.35 \pm 0.11$ \\
步骤，重新规划 & $3.4, 1.0$ & $4.9, 1.0$ & $18.0, 1.0$ & $5.9, 1.0$ & $7.6, 1.0$ & $12.6, 1.0$ \\
\midrule
对话（本文） &  &  &  &  &  &  \\
成功率 & $0.65 \pm 0.11$ & $0.93 \pm 0.06$ & $0.44 \pm 0.06$ & $0.75 \pm 0.10$ & $0.95 \pm 0.05$ & $0.80 \pm 0.08$ \\
步骤，重新规划 & $2.5, 3.1$ & $4.9, 1.3$ & $9.9, 3.5$ & $4.7, 2.0$ & $7.1, 1.0$ & $10.2, 1.7$ \\
\bottomrule
\end{tabular}
\end{table}

\subsection{大语言模型提出的三维路点的效果}

我们展示了大语言模型提出的任务空间路点的实用性。我们使用两个设计为具有高工作空间重叠的任务，即"包装杂货"和"移动绳子"，它们需要拾取和放置来完成任务。作为比较，我们定义了一个硬编码的路点路径，执行自上而下的拾取或放置，即总是在某个高度上方悬停夹爪后再拾取物体，并在移动和放置之前将物体提升到桌面上方的高度。我们抽取单个步骤的拾取或放置快照，并在最多300秒的规划时间预算下，使用比较的路点运行多臂运动规划。

大语言模型提出的路点在拾取子任务方面没有明显优势，但显著加速了放置规划，因为在这种情况下机械臂和桌面物体之间更容易发生碰撞。

\subsection{零样本适应任务变化}

利用大语言模型的零样本和少样本能力，RoCo 展示了对变化任务语义的强适应能力，这在传统上需要修改或重新编程系统，例如微调基于学习的策略。我们展示三个主要变化类别，全部使用 RoCoBench 中的"制作三明治"任务。

\begin{enumerate}
    \item \textbf{物体初始化：}食物物品的位置是随机的，我们展示对话智能体的推理对这种变化是稳健的。
    \item \textbf{任务目标：}智能体必须根据三明治食谱以正确顺序堆叠食物，并能相应地协调子任务策略。
    \item \textbf{机器人能力：}智能体能够交换关于各自可达范围内物品的信息，并相应地协调计划。
\end{enumerate}

\subsection{真实世界实验：人机协作}

我们在真实世界环境中验证 RoCo，其中一个机械臂与人类协作完成分类积木任务。我们 在修改后的 RoCo 上运行——只有机器人智能体由 GPT-4 控制，它与在部分任务工作空间中进行交互的人类用户讨论。对于感知，我们使用预训练的目标检测模型 OWL-ViT，从俯视 RGB-D 相机图像生成场景描述。任务约束人类只能将积木从杯子移动到桌子，然后机器人只能从桌子拾取积木放入木箱。我们评估两个主要变化类别：1）物体初始化，即每次运行的初始积木位置是随机的；2）任务顺序规范，即智能体被要求遵循固定顺序移动积木。我们还评估两种人类行为：第一种是 oracle 人类，纠正 OWL-ViT 引导的场景描述和机器人响应中的错误；第二种是不完美的人类，不对这些错误提供反馈。

我们对每种设置运行10次，表3报告了结果。我们报告有限轮次内的任务成功率，以及智能体成功完成一个事件所需的步骤数。我们注意到任务性能主要受 OWL-ViT 不正确的目标检测限制，这导致要么拾取错误的物体并导致失败，要么没有物体被拾取并导致步骤数增加。详见附录12.2了解真实世界实验的更多细节。

\begin{table}[htbp]
\centering
\caption{真实世界实验结果。我们报告成功次数和成功运行中的平均步数。}
\begin{tabular}{lcccc}
\toprule
 & \multicolumn{2}{c}{物体初始化} & \multicolumn{2}{c}{任务顺序} \\
\midrule
人类类型 & 成功率 & 步数 & 成功率 & 步数 \\
\midrule
Oracle 人类 & 9/10 & 5.3 & 8/10 & 5.5 \\
不完美人类 & 7/10 & 5.6 & 6/10 & 5.2 \\
\bottomrule
\end{tabular}
\end{table}

% 6. 多智能体表征与推理数据集
\section{多智能体表征与推理数据集}

除了我们的主要实验结果外，我们还整理了一个基于文本的数据集 RoCoBench-Text，用于评估大语言模型的智能体表征和任务推理能力。该数据集将大语言模型与多智能体协作中所需的能力对齐，无需机器人环境交互。它建立在我们对 RoCoBench 的评估数据基础上，并包含一系列更开放的问题和问题答案，超出了简单地找到下一个最佳动作计划。

\subsection{数据集概述}

该数据集包含是/否、多选或简短问答问题，涵盖一系列不同的推理能力（详见附录14）：

\textbf{自我知识}评估智能体在给定任务上下文中建立其身份的程度，分为两类：1）理解智能体自己的能力（如哪些物体/区域不可达）；2）记忆检索，即从过去的对话和动作中推断信息。

\textbf{沟通技巧}评估智能体有效交换信息和推动讨论达成一致计划的能力。问题要求大语言模型：1）选择对其他智能体问题的适当回复；2）选择对其他智能体的适当询问。

\textbf{适应能力}评估对上下文中未指定的意外情况的适应能力。我们使用 RoCoBench 任务的子集来设计意外事件，无论是关于任务状态（如缺失的物体）还是来自其他智能体的回复，并要求大语言模型智能体选择最佳回复。

详见下方示例问题：两个智能体一起制作三明治；其中一个智能体被告知夹爪损坏，必须推断三明治实际上可以在其侧不需要任何物品的情况下完成。

\subsection{评估结果}

\textbf{设置。}所有问题都设计为只有一个正确答案，因此我们测量每个类别的平均准确率。我们评估 GPT-4（OpenAI）、GPT-3.5-turbo（OpenAI）和 Claude-v1（Anthropic）。对于 GPT-4，我们使用标记有不同时间戳的两个模型，即03/14/2023和06/13/2023。结果总结于表4：我们观察到，在两个版本之间的小性能变化下，GPT-4 在所有类别中领先。我们注意到与完全准确仍有相当大的差距，希望这个数据集将有助于在未来工作中改进和评估语言模型。

\begin{table}[htbp]
\centering
\caption{多智能体大语言模型推理数据集的评估结果。我们测量每个测试类别的问答准确率，并比较四种不同模型的性能。}
\begin{tabular}{lccccc}
\toprule
模型 & 自我知识 & \multicolumn{2}{c}{沟通技巧} & 适应能力 \\
\midrule
 & 能力 & 记忆 & 询问 & 回复 \\
\midrule
GPT-4-0314 & $0.67 \pm 0.06$ & $0.84 \pm 0.06$ & $0.79 \pm 0.06$ & $0.83 \pm 0.04$ & $0.68 \pm 0.08$ \\
GPT-4-0613 & $0.68 \pm 0.06$ & $0.91 \pm 0.04$ & $0.57 \pm 0.08$ & $0.86 \pm 0.03$ & $0.71 \pm 0.08$ \\
GPT-3.5-turbo & $0.68 \pm 0.06$ & $0.59 \pm 0.07$ & $0.48 \pm 0.08$ & $0.30 \pm 0.05$ & $0.58 \pm 0.09$ \\
Claude-v1 & $0.37 \pm 0.06$ & $0.70 \pm 0.07$ & $0.55 \pm 0.08$ & $0.60 \pm 0.05$ & $0.65 \pm 0.09$ \\
\bottomrule
\end{tabular}
\end{table}

% 7. 局限性
\section{局限性}

\textbf{仿真中的 oracle 状态信息。}RoCo 假设感知（如目标检测、姿态估计和碰撞检测）是准确的。这一假设使我们的方法在无法获得完美感知的情况下容易失败：这反映在我们的真实世界实验中，预训练目标检测产生的错误可能导致规划错误。

\textbf{开环执行。}来自我们规划器的运动轨迹由机器人在开环方式下执行，并导致潜在错误。由于场景和动作描述中的抽象层，大语言模型无法识别或找到处理此类执行级错误的方法。

\textbf{大语言模型查询效率。}我们依赖查询预训练大语言模型来生成智能体对话中每个回复，这可能成本高昂且依赖于大语言模型的响应时间。对于动态或对速度敏感的任务，来自大语言模型查询的响应延迟是不可取的。

% 8. 相关工作
\section{相关工作}

\subsection{用于机器人学的大语言模型}

一系列早期工作使用大语言模型来选择技能原语并完成机器人任务，如 SayCan、Inner Monologue，它们与我们的工作类似，使用环境反馈在上下文中改进规划。后续工作利用大语言模型的代码生成能力以代码格式生成机器人策略，如 CaP、ProgGPT 和 Demo2Code；或为机器人执行生成长程序，如 TidyBot 和 Instruct2Act。与我们使用运动规划器相关的是 Text2Motion、AutoTAMP 和 LLM-GROP 研究将大语言模型与传统的任务与运动规划（Task and Motion Planning, TAMP）相结合。其他工作探索使用大语言模型促进人机协作、设计强化学习（Reinforcement Learning, RL）的奖励，以及在机器人任务中实时运动规划控制。虽然先前的工作使用单机器人设置和单线程大语言模型规划，但我们考虑多机器人设置可以实现更复杂的任务，并使用对话提示进行任务推理和协调。

\subsection{用于机器人学的多模态预训练}

大语言模型缺乏感知能力限制了其与机器人应用的结合。一种解决方案是使用视觉、语言和大规模机器人数据预训练新模型：多模态预训练的 PALM-E 实现感知和任务规划于单一模型；Interactive Language 和 DIAL 构建大规模语言注释的机器人轨迹数据集，用于训练可泛化的模仿策略。另一种解决方案是引入其他预训练模型，主要是视觉-语言模型（Vision-Language Model, VLM）如 CLIP。在 Socratic Models、Matcha 和 Kwon 等工作中，大语言模型被用来重复查询和综合来自其他模型的信息，以改进对环境的推理。虽然大多数使用零样本大语言模型和视觉-语言模型，但 CogLoop 也探索了微调适配层以更好地桥接不同的冻结模型。我们的工作利用仿真提取感知信息，我们的真实世界实验遵循先前的工作使用预训练目标检测模型生成场景描述。

\subsection{对话、辩论和大语言模型角色扮演}

在机器人学之外，大语言模型已被证明具有表征智能体意图和行为的能力，这使得在模拟环境（如基于文本的游戏和社交沙盒）中进行多智能体交互成为可能。最近的工作还表明，对话或辩论式提示可以提高大语言模型在人类对齐和广泛的目标导向任务上的性能。虽然先前的工作更专注于理解大语言模型行为或改进单一问题的解决方案，但我们的设置需要为每个智能体规划单独的动作，因此增加了讨论的复杂性和达成共识的难度。

\subsection{多机器人协作与运动规划}

多机器人操作研究有着悠久的历史。一类工作研究寻找无碰撞运动轨迹的低层问题。基于采样的方法是流行方法，其中提出了各种算法改进。最近的工作也探索了基于学习的方法作为替代。虽然我们的任务主要设置在静态场景中，但更多工作也研究了需要更复杂低层控制的更具挑战性的场景，如涉及动态物体或闭链运动学。另一类工作更专注于高层规划以分配和协调子任务，这与我们更相关。虽然大多数先前的工作为小规模任务定制系统（如家具组装），但我们强调我们方法对以少样本方式实现的各种任务的通用性。

% 9. 结论
\section{结论}

我们提出了 RoCo，这是一个利用大语言模型（Large Language Model, LLM）进行机器人协调和规划的多机器人协作新框架。我们引入了 RoCoBench，这是一个包含6个任务的多机器人操作基准测试，将向更广泛的研究社区开源。我们实验证明了我们方法的通用性和许多理想特性，如对变化任务语义的少样本适应，同时也指出了局限性和改进空间。我们的工作与最近探索利用大语言模型进行机器人应用力量的文献一致，并指明了未来研究这一方向的许多令人兴奋的机遇。

\textbf{致谢。}本工作部分由 NSF Award \#2143601、\#2037101 和 \#2132519 支持。我们要感谢 Google 提供 UR5 机器人硬件。本文中表达的观点和结论属于作者，不应被解释为必然代表赞助商的官方政策，无论是明示还是暗示。作者要感谢 Zeyi Liu、Zhenjia Xu、Huy Ha、Cheng Chi、Samir Gadre、Mengda Xu 和 Dominik Bauer 在整个项目中的富有成果的讨论，以及对稿件初稿的有益反馈。

% 参考文献
\newpage
\bibliographystyle{plain}
\bibliography{references}

% 附录
\appendix

\section{RoCoBench 详细文档}

RoCoBench 基于 MuJoCo 物理引擎构建。我们要感谢相关的开源工作，这些工作极大地帮助了 RoCoBench 任务的开发：DM-Control、Menagerie 和 Dasari 等人的 MuJoCo 物体资源。以下部分为6个模拟协作任务中的每一个提供详细文档。

\subsection{任务：清扫地面}

\textbf{任务描述。}两个机器人将扫帚和簸箕带到方块的相对两侧清扫它，然后拿着簸箕的机器人将方块倒入垃圾桶。

\textbf{智能体能力。}两个机器人站在桌子的相对两侧：1）配备 robotiq 夹爪的 UR5E（"Alice"）：拿着簸箕；2）Franka Panda（"Bob"）：拿着扫帚。

\textbf{观测空间。}1）方块位置：在桌子上；在簸箕内；在垃圾桶内；2）机器人状态：三维夹爪位置。

\textbf{可用技能。}1）MOVE [目标]：目标只能是方块；2）SWEEP [目标]：移动扫帚将目标推入簸箕；3）WAIT；4）DUMP：将簸箕倒在垃圾桶顶部。

\subsection{任务：制作三明治}

\textbf{任务描述。}两个机器人一起制作三明治，每个都可以接触到一组不同的食材。它们必须选择所需的物品并轮流按正确顺序堆叠。

\textbf{智能体能力。}两个机器人站在桌子的相对两侧：1）配备吸力夹爪的 UR5E（"Chad"）：只能到达右侧；2）配备吸力夹爪的人形机器人（"Dave"）：只能到达左侧。

\textbf{观测空间。}1）机器人自己的夹爪状态（空或拿着物体）；2）自己一侧桌上的食物和砧板上的食物。

\textbf{可用技能。}1）PICK [物体]；2）PUT [物体] on [目标]；WAIT。

\subsection{任务：分类方块}

\textbf{任务描述。}3个机器人将3个方块分类到对应的面板上。机器人必须保持在各自的可达范围内，并相互帮助将方块移近。

\textbf{智能体能力。}三个机器人各自负责桌子的一个区域：1）配备 robotiq 夹爪的 UR5E（"Alice"）：必须将蓝色方块放到 panel2，只能到达：panel1、panel2、panel3。2）Franka Panda（"Bob"）：必须将粉色多边形放到 panel4，只能到达：panel3、panel4、panel5。3）配备吸力夹爪的 UR5E（"Chad"）：必须将黄色梯形放到 panel6，只能到达：panel5、panel6、panel7。

\textbf{观测空间。}1）机器人自己的目标，2）每个方块的位置。

\textbf{可用技能。}1）PICK [物体] PLACE [panelX]；2）WAIT。

\subsection{任务：包装杂货}

\textbf{任务描述。}2个机器人将一组杂货从桌子包装到箱子里。物体位置接近，机器人必须协调路径以避免碰撞。

\textbf{智能体能力。}桌子两侧的两个机器人：1）配备 robotiq 夹爪的 UR5E（"Alice"）：可以拾取和放置桌上的任何物体；2）Franka Panda（"Bob"）：可以拾取和放置桌上的任何物体。

\textbf{观测空间。}1）机器人夹爪位置，2）每个物体的位置，3）箱子所有槽的位置。

\textbf{可用技能。}（必须包含任务空间路点）1）PICK [物体] PATH [路径]；2）PLACE [物体] [目标] PATH [路径]。

\subsection{任务：移动绳子}

\textbf{任务描述。}2个机器人一起将绳子抬起越过墙壁并放入凹槽。它们必须协调夹爪以避免碰撞。

\textbf{智能体能力。}桌子两侧的两个机器人：1）配备 robotiq 夹爪的 UR5E（"Alice"）：可以拾取和放置其可达范围内绳子的任何一端；2）Franka Panda（"Bob"）：可以拾取和放置其可达范围内绳子的任何一端。

\textbf{观测空间。}1）机器人夹爪位置，2）绳子前后端的位置；3）障碍墙角的位置；4）凹槽槽左右端的位置。

\textbf{可用技能。}（必须包含任务空间路点）1）PICK [物体] PATH [路径]；2）PLACE [物体] [目标] PATH [路径]。

\subsection{任务：整理橱柜}

\textbf{任务描述。}3个机器人，其中两个各持橱柜门一侧打开，而第三个机器人取出杯子并将其放到正确的杯垫上。

\textbf{智能体能力。}三个机器人，一个在桌子左侧，两个在右侧：1）配备 robotiq 夹爪的 UR5E（"Alice"）：站在左侧，只能到达左侧橱柜门；2）Franka Panda（"Bob"）：站在右侧，只能到达右侧橱柜门；3）配备吸力夹爪的 UR5E（"Chad"）：站在右侧，可以到达右侧橱柜门和橱柜内的杯子和马克杯。

\textbf{观测空间。}1）橱柜门把手位置；2）每个机器人可触及的物体，不知道其他机器人的可达范围。

\textbf{可用技能。}1）PICK [物体]；2）OPEN [门把手一侧]；3）WAIT；3）PLACE [物体] [目标]。

\section{大语言模型提示详细说明}

我们为对话中每个智能体的单独回复使用单独的查询调用，详见下方智能体提示的编辑示例。

\section{玩具示例：大语言模型用于三维路径规划}

我们使用 GPT-4 的 chatCompletion 模式。在每次评估运行中，我们随机采样一组新的障碍物和智能体的起点-终点位置。每次运行给予最多5次尝试：如果第一次尝试失败，则将先前计划的反馈附加到后续尝试的用户提示中，直到达到最大尝试次数。

\section{多智能体表征与推理数据集}

该数据集包含是/否、多选或简短问答问题，涵盖一系列不同的推理能力。

\subsection{自我知识问答}

\textbf{智能体能力。}57个问题。使用顺序运输模拟任务。

\textbf{记忆检索。}44个问题。使用制作三明治和清扫地面任务。

\subsection{有效沟通挑战}

\textbf{询问。}41个多选题，使用整理橱柜任务。

\textbf{回复。}96个问题。二元是/否问题。

\subsection{适应意外场景挑战}

31个多选题（A、B、C）。

\end{document}
